% !TEX root = ../discretemath.tex

\section{Моменты случайной величины}


\subsection{Основные определения}

\begin{Definition}{Моменты случайной величины}{}
	Пусть $\xi$ — случайная величина, $k\in\mathbb N$.
	\begin{itemize}
		\item \textbf{Начальный момент} порядка $k$:
		\[
		\mu'_k=\mathbb{E}[\xi^k].
		\]
		При $k=1$: $\mu'_1=\mathbb{E}[\xi]$ — математическое ожидание.
		\item \textbf{Центральный момент} порядка $k\ge 2$:
		\[
		\mu_k=\mathbb{E}\bigl[(\xi-\mathbb{E}[\xi])^k\bigr].
		\]
		По определению $\mu_1=0$, а $\mu_2=\mathbb{E}\bigl[(\xi-\mathbb{E}[\xi])^2\bigr]$ — дисперсия.
	\end{itemize}
\end{Definition}

\begin{Remark}{}{}
	Начальный момент $\mu'_k$ «привязан к нулю», центральный $\mu_k$ — «к среднему». Первый момент задаёт положение, второй центральный (дисперсия) — разброс; моменты высших порядков характеризуют форму распределения (асимметрию, остроту).
\end{Remark}

\subsection{Связь центральных и начальных моментов}

\begin{Theorem}{Выражение центрального момента через начальные}{central-via-initial}
	Для любого $k\ge 2$:
	\[
	\mu_k=\sum_{s=0}^{k}(-1)^s\binom{k}{s}\bigl(\mathbb{E}[\xi]\bigr)^s\cdot\mu'_{k-s}.
	\]
\end{Theorem}

\begin{proof}
	Раскроем степень по биному Ньютона и применим линейность матожидания:
	\begin{align*}
		\mu_k
		&=\mathbb{E}\bigl[(\xi-\mathbb{E}[\xi])^k\bigr]
		=\mathbb{E}\!\left[\sum_{s=0}^{k}(-1)^s\binom{k}{s}(\mathbb{E}[\xi])^s\xi^{k-s}\right]\\
		&=\sum_{s=0}^{k}(-1)^s\binom{k}{s}(\mathbb{E}[\xi])^s\,\mathbb{E}[\xi^{k-s}]
		=\sum_{s=0}^{k}(-1)^s\binom{k}{s}(\mathbb{E}[\xi])^s\,\mu'_{k-s},
	\end{align*}
	где $(\mathbb{E}[\xi])^s$ — константа, вынесена из-под матожидания, и $\mathbb{E}[\xi^{k-s}]=\mu'_{k-s}$.
\end{proof}

\begin{Corollary}{}{}
	Поскольку $\mathbb{E}[\xi]=\mu'_1$, формула принимает вид
	\[
	\mu_k=\sum_{s=0}^{k}(-1)^s\binom{k}{s}(\mu'_1)^s\,\mu'_{k-s}.
	\]
\end{Corollary}

\begin{Corollary}{}{}
	Каждый центральный момент $\mu_k$ выражается через начальные $\mu'_1,\dots,\mu'_k$, и наоборот: каждый начальный $\mu'_k$ выражается через $\mu'_1$ и центральные $\mu_2,\dots,\mu_k$.
\end{Corollary}

\begin{Remark}{}{}
	Это возможно потому, что наборы $\{1,\xi,\xi^2,\ldots,\xi^k\}$ и $\{1,(\xi-\mathbb E\xi),(\xi-\mathbb E\xi)^2,\ldots,(\xi-\mathbb E\xi)^k\}$ — базисы одной и той же линейной оболочки. Переход между ними — треугольное (унитреугольное) преобразование, потому обратимое.
\end{Remark}

\subsection{Единственность распределения по моментам}

\begin{Theorem}{Моменты определяют распределение}{}
	Пусть $\xi$ и $\beta$ — случайные величины над конечным вероятностным пространством $P$, и все их начальные моменты совпадают: $\mu'_k(\xi)=\mu'_k(\beta)$ для всех $k$. Тогда $\xi\overset{d}{=}\beta$ (одинаковое распределение).
\end{Theorem}

\begin{proof}
	Пусть $x_1,\dots,x_n$ — все возможные значения $\xi$ и $\beta$ (объединение). Обозначим $d_\xi(i)=P[\xi=x_i]$, $\sum_i d_\xi(i)=1$, и аналогично $d_\beta(i)=P[\beta=x_i]$.

	Рассмотрим матрицу Вандермонда
	\[
	M=\begin{pmatrix}
		1 & 1 & \cdots & 1\\
		x_1 & x_2 & \cdots & x_n\\
		x_1^2 & x_2^2 & \cdots & x_n^2\\
		\vdots & \vdots & \ddots & \vdots\\
		x_1^{n-1} & x_2^{n-1} & \cdots & x_n^{n-1}
	\end{pmatrix}.
	\]
	Тогда
	\[
	M\begin{pmatrix}d_\xi(1)\\\vdots\\d_\xi(n)\end{pmatrix}
	=\begin{pmatrix}\sum_i d_\xi(i)\\\sum_i x_i d_\xi(i)\\\vdots\\\sum_i x_i^{n-1}d_\xi(i)\end{pmatrix}
	=\begin{pmatrix}1\\\mu'_1(\xi)\\\vdots\\\mu'_{n-1}(\xi)\end{pmatrix}.
	\]
	Так как моменты совпадают, $M(d_\xi-d_\beta)=\mathbf 0$.

	\textbf{Невырожденность $M$.} Определитель $\det M$ — однородный многочлен от $x_1,\dots,x_n$ степени $\tfrac{n(n-1)}{2}$. Он обращается в нуль при $x_i=x_j$, поэтому каждый множитель $(x_i-x_j)$ делит $\det M$; так как
	\[
	\deg(\det M)=\deg\prod_{i<j}(x_i-x_j)=\frac{n(n-1)}{2},
	\]
	то $\det M=C\prod_{i<j}(x_i-x_j)$. Сравнивая коэффициент при мономе $x_1^0x_2^1\cdots x_n^{n-1}$ (он равен $1$ в обоих выражениях), получаем $C=1$, откуда
	\[
	\det M=\prod_{i<j}(x_i-x_j)\ne 0
	\]
	при попарно различных $x_i$.

	Значит, $M$ невырождена, и из $M(d_\xi-d_\beta)=\mathbf 0$ следует $d_\xi(i)=d_\beta(i)$ для всех $i$, то есть $\xi$ и $\beta$ имеют одинаковое распределение.
\end{proof}

\begin{Remark}{}{}
	Конечность пространства существенна: достаточно первых $n-1$ моментов (по числу значений). Для бесконечных распределений совпадение всех моментов в общем случае распределение \emph{не} определяет (проблема моментов).
\end{Remark}
